%%
%% This is file `sample-sigconf.tex',
%% generated with the docstrip utility.
%%
%% The original source files were:
%%
%% samples.dtx  (with options: `all,proceedings,bibtex,sigconf')
%% 
%% IMPORTANT NOTICE:
%% 
%% For the copyright see the source file.
%% 
%% Any modified versions of this file must be renamed
%% with new filenames distinct from sample-sigconf.tex.
%% 
%% For distribution of the original source see the terms
%% for copying and modification in the file samples.dtx.
%% 
%% This generated file may be distributed as long as the
%% original source files, as listed above, are part of the
%% same distribution. (The sources need not necessarily be
%% in the same archive or directory.)
%%
%%
%% Commands for TeXCount
%TC:macro \cite [option:text,text]
%TC:macro \citep [option:text,text]
%TC:macro \citet [option:text,text]
%TC:envir table 0 1
%TC:envir table* 0 1
%TC:envir tabular [ignore] word
%TC:envir displaymath 0 word
%TC:envir math 0 word
%TC:envir comment 0 0
%%
%% The first command in your LaTeX source must be the \documentclass
%% command.
%%
%% For submission and review of your manuscript please change the
%% command to \documentclass[manuscript, screen, review]{acmart}.
%%
%% When submitting camera ready or to TAPS, please change the command
%% to \documentclass[sigconf]{acmart} or whichever template is required
%% for your publication.
%%
%%
\documentclass[sigconf]{acmart}

\usepackage{fancyhdr}
% \usepackage{algorithmic}
\usepackage[linesnumbered,vlined,ruled]{algorithm2e}
\usepackage{algpseudocode}
\usepackage{graphicx}
\usepackage{textcomp}
\usepackage{xcolor}
\usepackage{bm}
\usepackage{booktabs}
\usepackage{multirow}
\usepackage{enumitem}
\usepackage{microtype}
% \usepackage{graphicx}
\usepackage{subfigure}
\usepackage{booktabs} % for professional tables
\usepackage{amsfonts} 
\let\Bbbk\relax
\usepackage{amssymb}
\usepackage{amsmath}
\usepackage{amsthm}
\usepackage{sansmath}
%\usepackage{MnSymbol,wasysym}
\usepackage{subfigure}
\usepackage{amsthm}
\usepackage{tcolorbox}
\usepackage{colortbl}
\usepackage{tikz}
\usepackage{url}
\usepackage{pifont}
% \usepackage{graphicx}
%\usepackage{subcaption}
\usepackage{makecell}
\usepackage{cleveref}
\usepackage{pgfplots}
\usepackage{setspace}
\setlength{\textfloatsep}{5pt}  % 图表与文本之间的间距
\setlength{\floatsep}{5pt}       % 图表之间的间距
\setlength{\intextsep}{5pt} 
\usepackage[dvipsnames]{xcolor}

% Remove ACM copyright and other information
\settopmatter{printacmref=false} % Disable ACM Reference Format
\renewcommand\footnotetextcopyrightpermission[1]{} % Remove copyright footnote


\newcommand{\CSV}{\textsc{CSV}\xspace}
\newcommand{\CSVU}{\textsc{UniCSV}\xspace}
\newcommand{\CSVS}{\textsc{SimCSV}\xspace}
\newcommand{\CSVUJ}{\textsc{UniJoin}\xspace}
\newcommand{\CSVSJ}{\textsc{SimJoin}\xspace}
\newcommand{\CSVNUJ}{\textsc{UniNeJoin}\xspace}
\newcommand{\CSVNSJ}{\textsc{SimNeJoin}\xspace}
\newcommand{\basefilter}{{Reference}\xspace}
\newcommand{\basejoin}{{CartJoin}\xspace}
\newcommand{\lotus}{{Lotus}\xspace}

\newcommand{\deepcolor}{\cellcolor[rgb]{0.8745, 0.8235, 0.7137}}
\newcommand{\lightcolor}{ \cellcolor[rgb]{0.9882,0.9373,0.8196}}

\crefformat{section}{\S#2#1#3} % see manual of cleveref, section 8.2.1
\crefformat{subsection}{\S#2#1#3}
\long\def\comment#1{}
\newcommand{\kfadd}[1]{\textcolor{blue}{#1}}
\newcommand{\PtoP}{{\textsc{P2P}}\xspace}
\newcommand{\MBTI}{\text{MBTI}\xspace}
%%
%% \BibTeX command to typeset BibTeX logo in the docs
\AtBeginDocument{%
  \providecommand\BibTeX{{%
    Bib\TeX}}}

\setcopyright{acmcopyright}
\copyrightyear{2026}
\acmYear{2026}
\acmDOI{XXXXXXX.XXXXXXX}

\acmConference[SIGMOD '26]{the 2026 International Conference on Management of Data}{May 31--June 5,
  2026}{Bengaluru, India}

\acmPrice{15.00}
\acmISBN{978-1-4503-XXXX-X/18/06}

\begin{document}

%%
%% The "title" command has an optional parameter,
%% allowing the author to define a "short title" to be used in page headers.
\title{Cost-effective Multi-way Merge Semantic ORDER BY \\ (Revision Plan)}

%%
%% The "author" command and its associated commands are used to define
%% the authors and their affiliations.
%% Of note is the shared affiliation of the first two authors, and the
%% "authornote" and "authornotemark" commands
%% used to denote shared contribution to the research.


% Core type and local formatting helpers.
\newcommand{\kw}[1]{{\ensuremath {\mathsf{#1}}}\xspace}
\newcommand{\kwnospace}[1]{{\ensuremath {\mathsf{#1}}}}
\newcommand{\etitle}[1]{\par\addvspace{1ex}\noindent\underline{\emph{#1}}}
\newcommand{\algocomment}[1]{{\footnotesize $\rhd$ \emph{#1}}}
\newcommand{\bigo}{\ensuremath{\mathcal{O}}}%
\newcommand{\tabincell}[2]{\begin{tabular}{@{}#1@{}}#2\end{tabular}}
\newcommand{\xmark}{\ding{55}}%

% Migrated passages use single-label lowercase section references, plus
% capitalized figure/table references. Avoid enabling cleveref globally:
% in this build it drops some existing amsmath labels under LaTeX 2026.
\providecommand{\cref}[1]{\S\ref{#1}}
\providecommand{\Cref}[1]{\autoref{#1}}

%%
%% By default, the full list of authors will be used in the page
%% headers. Often, this list is too long, and will overlap
%% other information printed in the page headers. This command allows
%% the author to define a more concise list
%% of authors' names for this purpose.
%\renewcommand{\shortauthors}{Trovato et al.}

%%
%% The abstract is a short summary of the work to be presented in the
%% article.


% \begin{CCSXML}
% <ccs2012>
%    <concept>
%        <concept_id>10002951.10003317.10003338.10003341</concept_id>
%        <concept_desc>Information systems~Language models</concept_desc>
%        <concept_significance>500</concept_significance>
%        </concept>
%    <concept>
%        <concept_id>10002951.10003317.10003347.10003349</concept_id>
%        <concept_desc>Information systems~Document filtering</concept_desc>
%        <concept_significance>500</concept_significance>
%        </concept>
%  </ccs2012>
% \end{CCSXML}

% \ccsdesc[500]{Information systems~Language models}
% \ccsdesc[500]{Information systems~Document filtering}

%%
%% Keywords. The author(s) should pick words that accurately describe
%% the work being presented. Separate the keywords with commas.
% \keywords{Large Language Model, Semantic Filter, Semantic Join}



% \received{20 February 2007}
% \received[revised]{12 March 2009}
% \received[accepted]{5 June 2009}

%%
%% This command processes the author and affiliation and title
%% information and builds the first part of the formatted document.
\maketitle
We thank all the reviewers for their invaluable and insightful comments. In the following, we provide the point-to-point response for each reviewer, where the revision plans are \revise{highlighted}.


\vspace{-2mm}
\section*{Response to Reviewer \# 1}
{\noindent \textbf{R1.D1}}: Questions w.r.t. embedding models and clustering. 
% D1. Key questions are (1) whether the embedding model can encode key characteristics of the data that can be broadly used in a large number of semantic filter conditions (2) whether clustering can group by key dimensions of the embedding that are relevant to the filter condition. On the second point above if clusters are organized by dimensions not relevant this can resort to (in the worst case) linear complexity.


{\noindent \underline{\sl Response}:} (1) Embeddings are abstract and vectorized representations of the original test, encapsulating hidden and rich semantic information for semantic understanding tasks
%\footnote{Hongliu Cao, Recent advances in text embedding: A Comprehensive Review of Top-Performing Methods on the MTEB Benchmark, CoRR, abs/2406.01607, 2024;}
\footnote{Nils Reimers and Iryna Gurevych, Sentence-BERT: Sentence Embeddings using Siamese BERT-Networks, Proc. of EMNLP/IJCNLP (1), pp. 3980–3990, 2019.}
%\footnote{Juri Opitz, Lucas Moeller, Andrianos Michail, Sebastian Padó, and Simon Clematide, Interpretable Text Embeddings and Text Similarity Explanation: A Survey, Proceedings of the 2025 Conference on Empirical Methods in Natural Language Processing, Suzhou, China, Association for Computational Linguistics, 2025. DOI: 10.18653/v1/2025.emnlp-main.1135.}
, and thereby can be used for processing general semantic filters. 
(2) As the embeddings are the representations of text in a hidden space, it is difficult to interpret them at the level of individual dimensions and grouped dimensions for specific predicates. 
It is worth noting that our approach does not fully rely on the offline clustering.  We also incorporate online recursive re-partition to adapt to a specific predicate. In the worst case, if the predicate is semantically `orthogonal' to the clustering, the algorithm resorts to linear complexity.

%Embeddings are a powerful representation technique, broadly adopted in numerous downstream tasks such as semantic search, classification, and clustering (TODO: citation here). Their strength lies in capturing latent semantic characteristics of the data, which makes them suitable for supporting a wide variety of filter conditions. At the same time, embeddings are inherently difficult to interpret at the level of individual dimensions, since each vector encodes entangled features rather than explicit attributes

%Our method groups embeddings and evaluates whether the grouped characteristics align with the filter condition through a voting mechanism. When the alignment is not satisfactory, we apply re-clustering to refine the grouping. This adaptive strategy ensures that clusters remain relevant to the semantic filter conditions, while the voting step provides a systematic check that prevents irrelevant dimensions from dominating. By iteratively refining clusters in this way, we maintain both semantic accuracy and computational efficiency, thereby avoiding the worst-case scenario of linear complexity in practice.

{\noindent \textbf{R1.D2}}: Other algorithm design options. 
% D2. There are several other design options I was wondering if authors explored and could shed insight their design choices (1) rather than sampling from each cluster use query embedding to subset clusters (2) aggregate embeddings of sampled data (for either positives or negative) to do similarity search (3) post LLM invocation in samples utilizing per cluster match rate (to determine difficulty/confidence for example and the choose small/large models accordingly (4) iterative sampling as opposed to fixed sampling ratio



{\noindent \underline{\sl Response}:}
Thanks for your suggestions. 
We conduct a preliminary testing for options (1)-(3) in Table~\ref{response:tab:options} in the response, and \revise{will add the complete results and analysis in the revision}.
 Specifically, for (1), we employ similarity search with the query embedding to identify representative samples, followed by a voting mechanism. 
 For (2), we introduce a new voting strategy that evaluates true/false outcomes based on the distance between the data embedding and the aggregate embeddings of sampled data (positive/negative). 
 For (3), Lotus employs this threshold learning strategy across the entire dataset. We apply it to each cluster.  As shown in Table~\ref{response:tab:options}, option (1) achieves comparable accuracy but at a higher cost of LLM calls, while options (2) and (3) show accuracy degradation.
 
 For (4), iterative sampling has the potential to achieve similar error bounds with smaller ratios, but it requires further refinement from theoretical proof to algorithmic design. 
 We will explore this as a future direction. 


 %We view this as an important avenue for further work.
\begin{table}[t]
    \centering
    \tiny
    \caption{Other Design Options}
    \vspace{-2ex}
    \label{response:tab:options}
    \setlength{\tabcolsep}{3pt}
    \begin{tabular}{c|ccc|ccc|ccc}
    \toprule
    \multirow{2}{*}{\bf Method} & \multicolumn{3}{c}{\textbf{RV}} & \multicolumn{3}{c}{\textbf{CB-Q2}} & \multicolumn{3}{c}{\textbf{AD-Q2}} \\
    \cmidrule(lr){2-4} \cmidrule(lr){5-7} \cmidrule(lr){8-10}
    & \textbf{Acc.} & \textbf{F1} &\textbf{LLM Call} & \textbf{Acc.} & \textbf{F1} & \textbf{LLM Call} & \textbf{Acc.} & \textbf{F1} & \textbf{LLM Call} \\
    \midrule
    (1)-UniVote & 0.9230 & 0.92734 & 404 & 0.7824 & 0.7575 & 7606 & 0.9193 & 0.8249 & 2474 \\
    (1)-SimVote & 0.9230 & 0.9273 & 404 & 0.7824 & 0.7577 & 7625 & 0.9193 & 0.8249 & 2473 \\
    (2) & 0.9329 & 0.9327 & 404 & 0.7337 & 0.7056 & 404 & 0.8215 & 0.6684 & 2008 \\
    (3) & 0.9527 & 0.9531 & 72239 & 0.6929 & 0.5922 & 13947 & 0.8989 & 0.7770 & 596814 \\
    Ours-UniVote & 0.9481 & 0.9336 & 404 & 0.7782 & 0.7601 & 7378 & 0.9139 & 0.8249 & 2008 \\
    Ours-SimVote & 0.9305 & 0.9335 & 404 & 0.7813 & 0.7555 & 7378 & 0.9193 & 0.8249 & 2008 \\
    \bottomrule
    \end{tabular}
\end{table}


{\noindent \textbf{R1.D3}}: Benchmarks are small.
% D3. Benchmarks are very small with 7 and 4 queries. It seems all of the benchmarks (except those for semantic join) also operate on a single text column. While technically this is a semantic filter but also in many ways it is a significantly reduced form of the operator (binary classification)
%Multiclass filter: can our method apply to that situation?

{\noindent \underline{\sl Response}:}
\revise{In the revision, we will test additional queries and datasets to further demonstrate the robustness of our approach.} Specifically, for the RV dataset, we plan to add two new queries: \textit{"The review explicitly recommends that others watch"} and \textit{"The review focuses on factual descriptions like technical details"}. For the CB dataset, we will add the query \textit{"The AboutMe focuses on factual identification details".} 
Moreover, we plan to incorporate two datasets that are not used in prior studies, i.e., a toxic comment classification and a fact-checking task, both in the domain of NLP. 
%Moreover, there are various binary classification tasks in the NLP field that have not been used in prior work. We will incorporate two additional datasets, one for toxic comment classification and another is fact checking.

% {\color{red} 
% %We will expand the current dataset by adding additional queries to strengthen the benchmarks (TODO: new queries)
% In our current benchmarks, we indeed use an NLP binary classification task i.e., the IMDB‑review dataset for semantic positive vs. negative, to evaluate the semantic filter operator.
% }



{\noindent \textbf{R1.D4}}: How to process multiple columns?
% D4. How are multiple columns considered? Other than embedding the whole row I see no consideration / exploitation of multiple columns independently where for example different cardinalities can be exploited. 

{\noindent \underline{\sl Response}:}
In this work, we follow the definition and test settings of existing studies, e.g., Lotus, concentrating on designing new algorithms with a performance guarantee for filters on a single column.
%In this work, we focus on designing an efficient algorithm with performance guarantee, thereby we currently consider semantic filter on a single column. 
Our approach can be extended to support filters on multiple columns by generating a fused semantics embedding for multiple columns. 
For example, concatenating the text of multiple columns together with the column names and feeding the long context into a pre-trained LM to produce the row-wise embedding, or indexing the embeddings for all columns by hierarchical clustering. Thanks for your question. We will continue to study the extension in this direction.

% \textbf{Original: }
% We currently focus on the basic form of the filter, which operates on a single column, as this provides a powerful foundation for further work. For multiple columns, we consider two situations:
% \begin{itemize}
%     \item When a filter condition involves multiple columns ({\color{green}similar to a semantic join}), we merge those columns and construct a single embedding for each row. This approach allows us to continue with the same clustering and voting framework.
%     \item When there are multiple independent filter conditions, the problem becomes a complex optimization across multiple operators. Since our work is designed to optimize complexity within a single operator, it is not our focus. We agree that exploiting column-specific cardinalities could be beneficial, and we view cardinality estimation in multi-operator scenarios as an important direction for further work.
% \end{itemize}

% \textbf{Add:}

% In our current work, we focus on a single-column situation, as this allows us to isolate and evaluate the core mechanism clearly. The approach can be directly extended to multi-column scenarios. For each individual column, encode each data point into embeddings and perform clustering separately. For each combination of columns, merge the data into a single column, encode it into embeddings, and then perform clustering. Typically, only semantic columns, i.e., those containing complex content such as text or links, require building an offline index, so the overall complexity remains manageable. In extreme multi-column settings, the extension method may become costly. We view this complex optimization problem as an important area for further work.

{\noindent \textbf{R1.D5}}:  If computing separate embeddings is considered.
% D5. Also it is not clear if authors considered computing separate embeddings (in addition to whole row). 


{\noindent \underline{\sl Response}:}
As the response in R1.D4, we can compute individual embeddings for each column and construct a hierarchical clustered index in the offline stage. It may also entail a substantial extension of the theoretical proof. 
In addition, constructing a benchmark containing multi-column semantic filters, with practical queries and reliable ground-truth, is also a challenging task. We will delve into the extension of multi-column filters as future work.
%As noted in D4, our current approach merges multiple columns into a single embedding for each row. {\color{green}Concatenating embeddings from different columns preserves the information contained in each column while allowing us to construct a unified representation for downstream clustering and voting}

% \textbf{Add:}
% As indicated in D4, we use both seperate embeddings and combine embeddings.


{\noindent \textbf{R1.D6}}: Incorporate auxiliary information into embeddings. 
% D6. More information would also be useful if table schema is incorporated into computing embeddings.

{\noindent \underline{\sl Response}:}
To incorporate the table schema into embeddings, we concatenate the column name to each value and encode the concatenated text into embeddings.
Table~\ref{response:tab:auxiliary} in the response shows the comparison results between using the new embeddings and our original performance, which demonstrates that the auxiliary column names impact the performance slightly. 


\begin{table}[t]
    \centering
    \tiny
    \caption{Incorporation Metadata into Embeddings}
    \vspace{-2ex}
    \label{response:tab:auxiliary}
    \setlength{\tabcolsep}{3pt}
    \begin{tabular}{c|ccc|ccc|ccc}
    \toprule
    \multirow{2}{*}{\bf Method} & \multicolumn{3}{c}{\textbf{RV}} & \multicolumn{3}{c}{\textbf{CB-Q2}} & \multicolumn{3}{c}{\textbf{AD-Q2}} \\
    \cmidrule(lr){2-4} \cmidrule(lr){5-7} \cmidrule(lr){8-10}
    & \textbf{Acc.} & \textbf{F1} &\textbf{LLM Call} & \textbf{Acc.} & \textbf{F1} & \textbf{LLM Call} & \textbf{Acc.} & \textbf{F1} & \textbf{LLM Call} \\
    \midrule
    UniVote & 0.9229 & 0.9272 & 404 & 9.7846 & 0.7612 & 7643 & 0.9243 & 0.8347 & 2008 \\
    SimVote & 0.9230 & 0.9273 & 404 & 0.7851 & 0.7617 & 7601 & 0.9242 & 0.8345 & 2008 \\
    Ours-UniVote & 0.9481 & 0.9336 & 404 & 0.7782 & 0.7601 & 7378 & 0.9139 & 0.8249 & 2008 \\
    Ours-SimVote & 0.9305 & 0.9335 & 404 & 0.7813 & 0.7555 & 7378 & 0.9193 & 0.8249 & 2008 \\
    \bottomrule
    \end{tabular}
\end{table}

{\noindent \textbf{R1.D7}}: Number of clusters hyperparameter.
% D7: Re: number of clusters hyper parameter. To me the increase in the number of LLM calls suggests that the cluster structure isn’t helping (aligned to query semantics) performance.

{\noindent \underline{\sl Response}:}
The performance of our approach is robust to hyperparameters such as sampling ratio and \# clusters. 
We clarify that increasing \# LLM calls will improve the confidence of voting within clusters, but it does not mean the cluster structure is helpless. Clusters delineate the groups for voting, while sampled LLM calls determine how many tuples participate in the voting. If there are a large number of LLM calls, their voting results may tie, and a re-partition is required for the remaining tuples. 
%Our experiments show that efficiency is generally better when the number of clusters is set to $4$ rather than $2$. Also, to ensure stability, we introduce a hyperparameter to constrain the minimum sample size in each cluster, defined as $\max (\text{min\_sample}, \text{sample ratio} \times \text{cluster size})$. As the number of clusters increases, cluster sizes become smaller, and the sample size is often determined by the minimum sample rather than the sample ratio. This explains why the required number of LLM calls increases with more clusters in our experiments. Considering the consistency of LLM, this safeguard is necessary to ensure consistent and reliable performance. 


{\noindent \textbf{{R1.D8}}}: Analysis of F1 for UNICSV/SIMCSV on CB-Q1. 
%D8: I think the F1 score for UNICSV and SIMCSV on CB-Q1 needs deeper dive. I do not quite understand authors point on structurally defined query. This also connects to other finding where authors suggest embedding quality had little impact on ultimate result query. These can suggest two things (1) any grouping could work (2) non-alignment of cluster / query. For further experimentation I would suggest authors to use the same dataset for different semantic filter conditions.



{\noindent \underline{\sl Response}:}
\revise{In the revision, we will test additional queries and datasets to further demonstrate the robustness of our approach.} The CB-Q1 predicate `contains a social link' is a predicate referring to a specific concept, like traditional SQL query `\kw{LIKE}'. It exhibits an extremely low selectivity (0.0328), and the offline clustering is difficult to capture its specific filtering condition.
%{\color{red}As R4.o3 for CB-Q1. In our experiments across diverse embedding models, the results show a small difference. This does not imply that any grouping would work. Rather, it reflects that the embedding quality isn't the bottleneck.}

% 现在回答的不是他问的呀？？？

\vspace{-2mm}
\section*{Response to Reviewer \# 2}
{\noindent \textbf{R2.O1}}: 

{\noindent \underline{\sl Response}:}


{\noindent \textbf{R2.O2}}: 

{\noindent \underline{\sl Response}:}


{\noindent \textbf{R2.O3}}: 

{\noindent \underline{\sl Response}:}



{\noindent \textbf{R3.O4}}: 

{\noindent \underline{\sl Response}:}



\vspace{-2mm}
\section*{Response to Reviewer \# 4}
{\noindent \textbf{R4.D1}}: 

{\noindent \underline{\sl Response}:}


{\noindent \textbf{R4.D2}}: 

{\noindent \underline{\sl Response}:}



{\noindent \textbf{R4.D3}}: 

{\noindent \underline{\sl Response}:}


{\noindent \textbf{R4.D4}}: 

{\noindent \underline{\sl Response}:}



{\noindent \textbf{R4.D5}}: 

{\noindent \underline{\sl Response}:}




{\noindent \textbf{R4.D6}}: 

{\noindent \underline{\sl Response}:}


{\noindent \textbf{R4.D7}}: 

{\noindent \underline{\sl Response}:}



{\noindent \textbf{R4.D8}}: 

{\noindent \underline{\sl Response}:}


{\noindent \textbf{R4.D9}}: 

{\noindent \underline{\sl Response}:}



{\noindent \textbf{R4.D10}}: 

{\noindent \underline{\sl Response}:}


\end{document}
\endinput
%%
%% End of file `sample-sigconf.tex'.
